% GENERATED FILE. Edit canonical lessons, then run npm run generate:books.
% canonical-source-hash: fnv1a64:2b7f93d5d3dd2e63
% canonical-lessons: ZH-C485-yali, ZH-C485-jingshen, ZH-C485-shenxin, ZH-C485-xintai, ZH-C485-jianya, ZH-R485-review-481-485, ZH-R485-review-final-485

\chapter{Body and Mind}
\label{ch:zh-body-and-mind}
\hlchaptermodality{485}

\begin{chapteropening}
\textbf{By the end of this chapter:} \emph{I can talk about stress, energy, body and mind, my frame of mind and relieving stress.}

Complete the last lesson of chapter 485: I can talk about stress, energy, body and mind, my frame of mind and relieving stress.
\end{chapteropening}

\section[yālì]{yālì — pressure; stress}

\label{lesson:ZH-C485-yali}



\begin{quote}
Before the new one: say the Mandarin for worn out, then the Mandarin for resistance to illness.
\end{quote}

\subsection*{You'll want to know: yālì}
\emph{yālì} — \textquotedblleft{}pressure; stress\textquotedblright{}.

\begin{quote}
\emph{wǒ zuìjìn yālì hěn dà} — I've been under a lot of stress lately.
\end{quote}

Its characters come later; for now it is a word for the ear.

\begin{cousinweb}
\emph{yālì} is built from \emph{yā} 'to press' and \zh{力} 'strength', both inherited from Middle Chinese, \emph{'aep} and \emph{lik} (Guangyun, in Baxter's transcription), which glosses the first 'to press down' and the second 'sinew'.
\end{cousinweb}

\begin{grammarlens}[title={What you've built}]
One more word for saying how it was.
\end{grammarlens}

\subsection*{Your turn}
\begin{itemize}
\raggedright
  \item \emph{Say it:} \emph{yālì}
  \item \emph{Say it:} \emph{yālì}, once more
  \item \emph{Say it:} \emph{wǒ zuìjìn yālì hěn dà} — I've been under a lot of stress lately.
  \item \emph{Recall:} say the Mandarin for worn out, then the Mandarin for resistance to illness, then say \emph{yālì} again
\end{itemize}

\begin{tcolorbox}[breakable,colback=teal!4,colframe=teal!35!black,title={Before you move on}]
What does yālì mean? (\textquotedblleft{}Pressure; stress\textquotedblright{}.) Say it once more.
\end{tcolorbox}
\section[jīngshen]{jīngshen — vigour; get-up-and-go}

\label{lesson:ZH-C485-jingshen}



\begin{quote}
Before the new one: say the Mandarin for resistance to illness, then the Mandarin for stress and pressure.
\end{quote}

\subsection*{You'll want to know: jīngshen}
\emph{jīngshen} — \textquotedblleft{}vigour; get-up-and-go\textquotedblright{}.

\begin{quote}
\emph{hē le kāfēi yǐhòu wǒ gǎnjué hěn yǒu jīngshen} — After a coffee, I feel full of energy.
\end{quote}

Its characters come later; for now it is a word for the ear.

\begin{cousinweb}
\emph{jīngshen} is built from \emph{jīng} 'refined' and \emph{shén} 'spirit', both inherited from Middle Chinese, \emph{tsjeng} and \emph{zyin} (Guangyun, in Baxter's transcription), which glosses the first 'bright' and 'pure' and the second 'divine'. In this sense the second is said light.
\end{cousinweb}

\begin{grammarlens}[title={What you've built}]
One more word for saying how it was.
\end{grammarlens}

\subsection*{Your turn}
\begin{itemize}
\raggedright
  \item \emph{Say it:} \emph{jīngshen}
  \item \emph{Say it:} \emph{jīngshen}, once more
  \item \emph{Say it:} \emph{hē le kāfēi yǐhòu wǒ gǎnjué hěn yǒu jīngshen} — After a coffee, I feel full of energy.
  \item \emph{Recall:} say the Mandarin for resistance to illness, then the Mandarin for stress and pressure, then say \emph{jīngshen} again
\end{itemize}

\begin{tcolorbox}[breakable,colback=teal!4,colframe=teal!35!black,title={Before you move on}]
What does jīngshen mean? (\textquotedblleft{}Vigour; get-up-and-go\textquotedblright{}.) Say it once more.
\end{tcolorbox}
\section[\texorpdfstring{\zh{身心} (shēnxīn)}{身心 (shēnxīn)}]{\zh{身心} (shēnxīn) — body and mind}

\label{lesson:ZH-C485-shenxin}



\begin{quote}
Before the new one: say the Mandarin for stress and pressure, then the Mandarin for get-up-and-go.
\end{quote}

\subsection*{You'll want to know: \zh{身心}}
\textbf{\zh{身心}} — \emph{shēnxīn} — \textquotedblleft{}body and mind\textquotedblright{}.

\begin{quote}
\textbf{\zh{眼前最要紧的是身心都好}} \emph{yǎnqián zuì yàojǐn de shì shēnxīn dōu hǎo} — For now, what matters most is being well in body and mind.
\end{quote}

\begin{cousinweb}
\textbf{\zh{身心}} (\emph{shēnxīn}) is built from \zh{身} 'body' and \zh{心} 'heart', both inherited from Middle Chinese, \emph{syin} and \emph{sim} (Guangyun, in Baxter's transcription), which glosses the first 'oneself'.
\end{cousinweb}

\begin{grammarlens}[title={What you've built}]
One more word for saying how it was.
\end{grammarlens}

\subsection*{Your turn}
\begin{itemize}
\raggedright
  \item \emph{Say it:} \emph{shēnxīn}
  \item \emph{Say it:} \emph{shēnxīn}, once more
  \item \emph{Say it:} \emph{yǎnqián zuì yàojǐn de shì shēnxīn dōu hǎo} — For now, what matters most is being well in body and mind.
  \item \emph{Recall:} say the Mandarin for stress and pressure, then the Mandarin for get-up-and-go, then say \emph{shēnxīn} again
\end{itemize}

\begin{tcolorbox}[breakable,colback=teal!4,colframe=teal!35!black,title={Before you move on}]
What does \zh{身心} mean? (\textquotedblleft{}Body and mind\textquotedblright{}.) Say it once more.
\end{tcolorbox}
\section[xīntài]{xīntài — mindset; attitude of mind}

\label{lesson:ZH-C485-xintai}



\begin{quote}
Before the new one: say the Mandarin for get-up-and-go, then the Mandarin for body and mind.
\end{quote}

\subsection*{You'll want to know: xīntài}
\emph{xīntài} — \textquotedblleft{}mindset; attitude of mind\textquotedblright{}.

\begin{quote}
\emph{kǎoshì de shíhou xīntài zuì zhòngyào} — In an exam, your mindset matters most.
\end{quote}

Its characters come later; for now it is a word for the ear.

\begin{cousinweb}
\emph{xīntài} is built from \zh{心} 'heart' and \emph{tài} 'a manner', both inherited from Middle Chinese, \emph{sim} and \emph{thojH} (Guangyun, in Baxter's transcription), which glosses the second 'bearing and manner'.
\end{cousinweb}

\begin{grammarlens}[title={What you've built}]
One more word for saying how it was.
\end{grammarlens}

\subsection*{Your turn}
\begin{itemize}
\raggedright
  \item \emph{Say it:} \emph{xīntài}
  \item \emph{Say it:} \emph{xīntài}, once more
  \item \emph{Say it:} \emph{kǎoshì de shíhou xīntài zuì zhòngyào} — In an exam, your mindset matters most.
  \item \emph{Recall:} say the Mandarin for get-up-and-go, then the Mandarin for body and mind, then say \emph{xīntài} again
\end{itemize}

\begin{tcolorbox}[breakable,colback=teal!4,colframe=teal!35!black,title={Before you move on}]
What does xīntài mean? (\textquotedblleft{}Mindset; attitude of mind\textquotedblright{}.) Say it once more.
\end{tcolorbox}
\section[jiǎnyā]{jiǎnyā — to relieve stress}

\label{lesson:ZH-C485-jianya}



\begin{quote}
Before the new one: say the Mandarin for body and mind, then the Mandarin for a mindset.
\end{quote}

\subsection*{You'll want to know: jiǎnyā}
\emph{jiǎnyā} — \textquotedblleft{}to relieve stress\textquotedblright{}.

\begin{quote}
\emph{tīng yīnyuè kěyǐ jiǎnyā} — Listening to music helps relieve stress.
\end{quote}

Its characters come later; for now it is a word for the ear.

\begin{cousinweb}
\emph{jiǎnyā} is built from \emph{jiǎn} 'to lessen' and \emph{yā} 'to press', both inherited from Middle Chinese, \emph{keamX} and \emph{'aep} (Guangyun, in Baxter's transcription). The dictionary lists the first under two readings and glosses \emph{keamX} 'to reduce'; it glosses the second 'to press down'.
\end{cousinweb}

\begin{grammarlens}[title={What you've built}]
One more word for saying how it was.
\end{grammarlens}

\subsection*{Your turn}
\begin{itemize}
\raggedright
  \item \emph{Say it:} \emph{jiǎnyā}
  \item \emph{Say it:} \emph{jiǎnyā}, once more
  \item \emph{Say it:} \emph{tīng yīnyuè kěyǐ jiǎnyā} — Listening to music helps relieve stress.
  \item \emph{Recall:} say the Mandarin for body and mind, then the Mandarin for a mindset, then say \emph{jiǎnyā} again
\end{itemize}

\begin{tcolorbox}[breakable,colback=teal!4,colframe=teal!35!black,title={Before you move on}]
What does jiǎnyā mean? (\textquotedblleft{}To relieve stress\textquotedblright{}.) Say it once more.
\end{tcolorbox}
\section[(dialogue)]{Looking back: from Is It Serious? to Body and Mind}

\label{lesson:ZH-R485-review-481-485}



\begin{quote}
Say the Mandarin for serious or urgent and how an illness is going.
\end{quote}

\subsection*{Your turn}
\begin{itemize}
\raggedright
  \item \emph{Say it:} serious or urgent, how an illness is going, the cause of an illness, to become infected, a virus
  \item \emph{Say it:} the emergency department, to give medical treatment, a pill, to disinfect, to care for the sick
  \item \emph{Say it:} to start to improve, to get well again, nourishment, looking after one's health, chronic
  \item \emph{Say it:} to guard against, to give up drinking, to get into a habit, worn out, resistance to illness
  \item \emph{Say it:} stress and pressure, get-up-and-go, body and mind, a mindset, to relieve stress
  \item \emph{Read it:} \zh{紧} — each one aloud, then a word that holds it
\end{itemize}

\begin{tcolorbox}[breakable,colback=teal!4,colframe=teal!35!black,title={Before you move on}]
Serious or urgent? (\textbf{\zh{要紧}}, \emph{yàojǐn}.) How an illness is going? (\textbf{\zh{病情}}, \emph{bìngqíng}.) The cause of an illness? (\textbf{\zh{病因}}, \emph{bìngyīn}.) To become infected? (\emph{gǎnrǎn}.) A virus? (\emph{bìngdú}.) The emergency department? (\emph{jízhěn}.) To give medical treatment? (\emph{zhìliáo}.) A pill? (\emph{yàopiàn}.) To disinfect? (\emph{xiāodú}.) To care for the sick? (\textbf{\zh{护理}}, \emph{hùlǐ}.) To start to improve? (\emph{hǎozhuǎn}.) To get well again? (\emph{kāngfù}.) Nourishment? (\emph{yíngyǎng}.) Looking after one's health? (\emph{bǎojiàn}.) Chronic? (\emph{mànxìng}.) To guard against? (\emph{yùfáng}.) To give up drinking? (\emph{jièjiǔ}.) To get into a habit? (\emph{yǎngchéng}.) Worn out? (\emph{píláo}.) Resistance to illness? (\emph{dǐkànglì}.) Stress and pressure? (\emph{yālì}.) Get-up-and-go? (\emph{jīngshen}.) Body and mind? (\textbf{\zh{身心}}, \emph{shēnxīn}.) A mindset? (\emph{xīntài}.) To relieve stress? (\emph{jiǎnyā}.)
\end{tcolorbox}
\section[(dialogue)]{Body and Mind — once more}

\label{lesson:ZH-R485-review-final-485}



\begin{quote}
Say the Mandarin for stress and pressure and get-up-and-go.
\end{quote}

\subsection*{Your turn}
\begin{itemize}
\raggedright
  \item \emph{Say it:} stress and pressure, get-up-and-go, body and mind, a mindset, to relieve stress
\end{itemize}

\begin{tcolorbox}[breakable,colback=teal!4,colframe=teal!35!black,title={Before you move on}]
Stress and pressure? (\emph{yālì}.) Get-up-and-go? (\emph{jīngshen}.) Body and mind? (\textbf{\zh{身心}}, \emph{shēnxīn}.) A mindset? (\emph{xīntài}.) To relieve stress? (\emph{jiǎnyā}.)
\end{tcolorbox}
